\documentclass{article}

\usepackage[polish]{babel}
\usepackage{polski}
\usepackage[utf8]{inputenc}
\usepackage[T1]{fontenc}
 
\usepackage[margin=1.5in]{geometry} 

\usepackage{color} 
\usepackage{amsmath}
\numberwithin{equation}{subsection}

\usepackage{amssymb}
\usepackage{amsfonts}                                                                   
\usepackage{graphicx}                                                             
\usepackage{booktabs}
\usepackage{amsthm}
\usepackage{pdfpages}
\usepackage{hyperref}
\usepackage{etoolbox}

\makeatletter
\newenvironment{definition}[1]{%
    \trivlist
    \item[\hskip\labelsep\textbf{Definition. #1.}]
    \ignorespaces
}{%
    \endtrivlist
}
\newenvironment{fact}[1]{%
    \trivlist
    \item[\hskip\labelsep\textbf{Fact. #1.}]
    \ignorespaces
}{%
    \endtrivlist
}
\newenvironment{theorem}[1]{%
    \trivlist
    \item[\hskip\labelsep\textbf{Theorem. #1.}]
    \ignorespaces
}{%
    \endtrivlist
}
\newenvironment{information}[1]{%
    \trivlist
    \item[\hskip\labelsep\textbf{Information. #1.}]
    \ignorespaces
}{%
    \endtrivlist
}
\newenvironment{identities}[1]{%
    \trivlist
    \item[\hskip\labelsep\textbf{Identities. #1.}]
    \ignorespaces
}{%
    \endtrivlist
}
\makeatother

\title{Kryptografia - Zadanie 2.3}  
\author{Rafał Włodarczyk}
\date{INA 6, 2026}

\begin{document}

\maketitle

\section{Wprowadzenie}

Obliczenia prowadzone są w grupie cyklicznej $G=\langle g \rangle$, 
gdzie $g$ jest generatorem grupy, a $q=\text{ord}(G)$ jest rzędem grupy. 
Osoba podpisująca posiada parę $(a, A)$ - $a$ jest kluczem prywatnym, 
$g^a=A$ jest kluczem publicznym. Zakładamy, że logarytm dyskretny w grupie $G$ jest trudny do obliczenia,
tj. nie istnieje efektywny algorytm który przy danych $G,g,A$ wyznacza $a : g^a=A$.
Wszystkie działania skalarne są wykonywane w obrębie grupy $\mathbb{Z}_q$.

\section{Zadanie 2.3}

W tym zadaniu funkcja skrótu oznaczona jest wzorem $h=mR$, gdzie $m$ - wiadomość oraz $R \in G$.
Mamy do czynienia z następującym schematem podpisu:
\begin{enumerate}
    \item Podpisujący losuje skalar $r\in \mathbb{Z}_q$ i oblicza $R=g^r$.
    \item Podpisujący wyznacza $h=mR$.
    \item Podpisujący oblicza $s = [r + ah] \mod q$ i zwraca $\sigma=(R,s)$ - podpis wiadomości $m$.
    \item Weryfikator wyznacza $h=mR$, a następnie weryfikuje podpis za pomocą warunku:
    \begin{align*}
        g^s = RA^h
    \end{align*}
\end{enumerate}

\section{Podpunkt 1}

\subsection{Poprawność Algebraiczna}

W nagłówku listy 2 wskazane jest dowiedzenie m.in poprawności algebraicznej u uczciwego podpisującego.
Zakładamy że posiada on klucz $a$, wybiera $r$, liczy $R=g^r$, wyznacza $h=mR, s = r + ah \mod q$ i zwraca $\sigma=(R,s)$.\\

\noindent
Pokażmy, że taki podpis przechodzi weryfikację:
\begin{align*}
    P = RA^h = g^r \cdot g^{ah} = g^{r + ah} = g^s = L_{\square}
\end{align*}

\noindent
Schemat jest algebraicznie poprawny.

\subsection{Polecenie}

Wyjaśnij, w jakiej domenie iloczyn $m\cdot R$ miałby być liczony.

\subsection{Rozwiązanie}

Wiemy, że $R$ jest elementem grupy $G$, a $m_{\text{original}}$ jest wiadomością o nieokreślonej długości. 
Wynikiem działania $h=m\cdot R$ musi być element grupy $\mathbb{Z}_q$ - ponieważ $h$ jest elementem wykładnika (w weryfikacji $g^s=RA^h$). 
Naturalnym przedstawieniem $mR$ w $\mathbb{Z}_q$ jest reprezentacja wiadomości $m$ 
przez liczbę całkowitą $m_{\text{int}}$ oraz $R$ rzutowanego na element grupy $\mathbb{Z}_q$ poprzez funkcję $f: G \to \mathbb{Z}_q$ (można np. przyjąć
porządek z generatora $g$), oraz następnie zredukowanie tych wartości modulo $q$:
\begin{align*}
    h = \left((m_{\text{int}}) \cdot (f(R) \bmod q)\right) \mod q
\end{align*}
\noindent
W bezpiecznym schemacie funkcja $h$ brałaby obie wartości $m \in M$ oraz $R\in G$ w taki sposób, że
$h : M \times G \rightarrow \mathbb{Z}_q$ oraz charakteryzowałaby się cechami bezpiecznej funkcji skrótu,
np. $h=H(m||R)$, gdzie $H$ jest bezpieczną funkcją skrótu.

\section{Podpunkt 2}

\subsection{Polecenie}

Sprawdź, czy takie związanie wiadomości z $R$ daje odporność na fałszerstwo, czy tylko pozór związania.

\subsection{Rozwiązanie}

Związanie $h=mR$ nie daje odporności na fałszerstwo, a jedynie pozór związania, umożliwiając atak \texttt{Existential Forgery}.\\

\noindent
Przyjmujemy, że fałszerstwo to zdolność do znalezienia \textbf{dowolnej wiadomości} $m$
dla której istnieje sposób wyznaczenia $(R,s)$ przechodzący weryfikację, 
a uzyskany \textbf{bez znajomości klucza prywatnego} $a$.
Niech $R\in G$ jest przedstawiony następująco: 
\begin{align*}
    R=g^x A^y \text{ dla pewnych } x,y \in \mathbb{Z}_q
\end{align*}
Zobaczmy warunek weryfikacji:
\begin{align*}
    g^s = RA^h = g^x A^y A^h = g^x A^{y+h}
\end{align*}
Chcemy zredukować wykładnik $A$ do postaci $A^0=1 (\text{mod } q)$, aby pozbyć się części wymagającej klucza prywatnego.
\begin{align*}
    y+h \equiv 0 \mod q \\
    y \equiv -h \mod q \\
    y \equiv -mR \mod q \\
    m \equiv -yR^{-1} \mod q
\end{align*}
Przyjmujemy, że wyznaczenie odwrotności $R^{-1}$ jest w grupie $G$ łatwe, np. na mocy rozszerzonego Algorytmu Euklidesa.
Wtedy $s$ jest dany jako $s = x \pmod q$.\\

\noindent
Zobaczmy, że dla tak spreparowanej trójki $(m,R,s)=(-yR^{-1},g^xA^y,x)$ warunek weryfikacji jest spełniony:
\begin{align*}
    P &= RA^h  \\
      &= g^xA^y \cdot A^{-mR} \\
      &= g^x A^{y+((-yR^{-1})\cdot R)} \\
      &= g^x A^{y+(-y\cdot(R^{-1}R))} \\
      &= g^x A^{y-y} \\
      &= g^x A^0 \\
      &= g^x \\
      &= g^s = L_{\square}
\end{align*}

\noindent
Związanie wiadomości z $R$ poprzez funkcję $h=mR$ daje pozór związania, ponieważ $h$ jest w istocie funkcją liniową od $R$.
Zaburza to jedną z kluczowych własności funkcji skrótu $h$, która powinna być trudna do odwrócenia.

\section{Podpunkt 3}

\subsection{Polecenie}

Zbadaj, czy można dobrać podpis dla wybranej wiadomości bez znajomości sekretu $a$.

\subsection{Rozwiązanie}

\texttt{Universal Forgery Attack} jest w tym wypadku niemożliwy.\\

\noindent
Pomimo braku odporności na fałszerstwo, nie można w prosty sposób dobrać podpisu dla wybranej wiadomości bez znajomości sekretu $a$.
Zobaczmy, że jeśli $m$ jest już wybrane to po stronie weryfikatora ustalona jest zależność $h=mR$, a zatem musi zajść weryfikacja:
\begin{align*}
    g^s = RA^h = Rg^{amR}
\end{align*} 
Co jest równoznaczne z $g^s = g^{\log_g(R) + amR}$, a zatem $s = \log_g(R) + amR \pmod q$.\\

\noindent
Wobec tego musielibyśmy wykorzystać prawdziwą parę $(r,R)$ z relacją $g^r = R$, a to jest niemożliwe z ustalonym $m$ i
bez znajomości $a$, ponieważ $R$ jest elementem grupy $G$ i jego logarytm dyskretny jest trudny do obliczenia. W 
podpunkcie drugim udało się dobrać podpis, ponieważ akceptowaliśmy arbitralną z punktu widzenia weryfikatora
wiadomość $m$. 

\section{Podpunkt 4}

\subsection{Polecenie}

Porównaj ten schemat z zadaniem 2: czy problem jest zasadniczo ten sam, czy inny?

\subsection{Rozwiązanie}

Problem jest zasadniczo ten sam.\\

\noindent
W zadaniu drugim prezentowany schemat to $h=m + R$ zamiast $h=mR$.
Schematy są błędne ze względu na niepoprawne własności funkcji $h$, która nie jest bezpieczną funkcją skrótu.
W obu przypadkach funkcja $h$ jest złożona z prostych operacji algebraicznych, przez co jest odwracalna.

\noindent
Możemy w obu przypadkach wykonać atak \texttt{Existential Forgery} (tj. w podpunkcie 2):
\begin{itemize}
    \item Dla $h=m+R$: $m = -y-R \pmod q$
    \item Dla $h=mR$: $m= -yR^{-1} \pmod q$
\end{itemize}

\section{Wnioski}

Pomimo że schemat jest algebraicznie poprawny, to jednak jest niebezpieczny \\
- podatny na atak\\ \texttt{Existential Forgery}, czyli
dobranie takiej trójki $(m,R,s)$ (bez znajomości klucza prywatnego $a$), która pozytywnie przechodzi weryfikację.\\

\noindent
Schemat jest podatny na atak, ale nie pozwala na wykradnięcie klucza prywatnego $a$ (poza sytuacją, gdy
podpisujący powtórzy $R$ dla różnych wiadomości), wobec tego nie jest najgorszym schematem podpisu, wciąż jednak jest niebezpieczny.\\

\noindent
Kluczowa dla bezpieczeństwa schematu podpisu jest bezpieczna funkcja skrótu charakteryzująca
się odpornością na kolizję oraz trudnością odwracania - uwydatnione w zadaniu,
ale również m.in. determinizmem. 

\end{document}
